\chapter{MySQL / MariaDB}

\section{InnoDB vs MyISAM}

InnoDB es el motor default desde MySQL 5.5. MyISAM es legado --- sin transacciones, sin FK.

\begin{itemize}
  \item \textbf{InnoDB:} ACID, FK, row-level locking, MVCC, crash recovery
  \item \textbf{MyISAM:} más rápido en lecturas puras, sin transacciones, table-level lock
  \item Siempre usa InnoDB en producción. MyISAM solo para tablas de solo lectura con volumen extremo.
\end{itemize}

\section{Diferencias clave con PostgreSQL}

\begin{itemize}
  \item \texttt{AUTO\_INCREMENT} vs \texttt{SERIAL} (PG)
  \item No tiene \texttt{RETURNING} en INSERT/UPDATE
  \item \texttt{GROUP BY} más permisivo (columnas no agregadas sin ONLY\_FULL\_GROUP\_BY)
  \item \texttt{SHOW PROCESSLIST}, \texttt{SHOW STATUS} para monitoreo
  \item Replicación master-slave nativa (binlog)
\end{itemize}

\section{Particionamiento}

Divide tabla en particiones físicas. Tipos: RANGE, LIST, HASH, KEY. Mejora rendimiento en tablas grandes cuando las queries filtran por la columna de partición.

\begin{lstlisting}[caption={Sintaxis MySQL --- diferencias con PostgreSQL}]
-- AUTO_INCREMENT (MySQL) vs SERIAL (PostgreSQL)
CREATE TABLE productos (
    id      INT AUTO_INCREMENT PRIMARY KEY,
    nombre  VARCHAR(100) NOT NULL,
    precio  DECIMAL(10,2)
) ENGINE=InnoDB DEFAULT CHARSET=utf8mb4;

-- UPSERT en MySQL
INSERT INTO clientes (email, nombre, visitas)
VALUES ('ana@ej.com', 'Ana', 1)
ON DUPLICATE KEY UPDATE visitas = visitas + 1;

-- Ver motor de cada tabla
SELECT table_name, engine
FROM   information_schema.tables
WHERE  table_schema = 'mi_base';
\end{lstlisting}

\begin{lstlisting}[caption={Particionamiento por RANGE en MySQL}]
CREATE TABLE ventas (
    id      INT AUTO_INCREMENT,
    fecha   DATE NOT NULL,
    monto   DECIMAL(10,2),
    PRIMARY KEY (id, fecha)
)
PARTITION BY RANGE (YEAR(fecha)) (
    PARTITION p2022 VALUES LESS THAN (2023),
    PARTITION p2023 VALUES LESS THAN (2024),
    PARTITION p2024 VALUES LESS THAN (2025),
    PARTITION pfuture VALUES LESS THAN MAXVALUE
);

-- Partition pruning automático
EXPLAIN SELECT * FROM ventas
WHERE fecha BETWEEN '2023-01-01' AND '2023-12-31';
\end{lstlisting}

\begin{lstlisting}[caption={Monitoreo y replicación básica}]
-- Ver procesos activos
SHOW PROCESSLIST;

-- Ver estado del servidor
SHOW STATUS LIKE 'Threads%';
SHOW STATUS LIKE 'Innodb_row_lock%';

-- Estado de replicación (en esclavo)
SHOW SLAVE STATUS\G

-- Variables de configuración
SHOW VARIABLES LIKE 'innodb_buffer_pool_size';
SHOW VARIABLES LIKE 'max_connections';
\end{lstlisting}

\begin{entrevistabox}
\begin{enumerate}
  \item ¿Cuáles son las diferencias principales entre InnoDB y MyISAM?
  \item ¿Cómo implementarías un UPSERT en MySQL?
  \item ¿Qué es el partition pruning y cuándo beneficia al rendimiento?
\end{enumerate}
\end{entrevistabox}
